\section{GRPO Group Training Details}
\label{app:grpo_groups}

\begin{figure*}[h]
\centering
\includegraphics[width=\textwidth]{fig_grpo_pipeline.pdf}
\caption{Group-wise offline GRPO pipeline. Call types are partitioned into
four groups processed sequentially. Each group undergoes SFT on Opus
golden samples followed by offline GRPO with importance sampling. The
merged LoRA checkpoint from each group becomes the baseline for the next,
preventing cross-modality gradient interference.}
\label{fig:grpo_pipeline}
\end{figure*}

The debate pipeline produces four heterogeneous call types---tactic
selection, skeleton/evidence construction, speech generation, and
cross-examination dialogue---that differ dramatically in output length,
vocabulary, and evaluation criteria. Training a single GRPO objective
over all call types conflates these modalities and produces unstable
gradients. We instead partition calls into four groups (A--D), training
each sequentially with group-specific hyperparameters.

\subsection{Group Configurations}

\begin{table}[h]
\centering
\small
\begin{tabular}{lllrrrr}
\toprule
\textbf{Group} & \textbf{Call Types} & \textbf{Samples} & \textbf{Max Tok} & \textbf{Batch} & \textbf{Epochs} & \textbf{Val Acc} \\
\midrule
A & \textsc{Tactic\_Select} & 2,887 & 5,000 & 3--4 & 5 & 81.0\% \\
B & \textsc{Skeleton} + \textsc{Evidence} & 4,157 & 16,000 & 3 & 5 & 73.3\% \\
C & \textsc{Speech\_Generate} & 3,274 & 10,500 & 4 & 8 & 72.3\% \\
D & \textsc{CX\_Dialogue} & 5,003 & 7,000 & 4 & 4 & 87.5\% \\
\bottomrule
\end{tabular}
\caption{GRPO group configurations. Max tokens reflects the longest
completions in each group. Validation accuracy is measured on held-out
debates from the same iteration. Group~D achieves the highest validation
accuracy because CX responses have shorter, more constrained outputs.}
\label{tab:grpo_groups}
\end{table}

\paragraph{Why these groupings.}
Group~A (tactic selection) has the shortest outputs ($\sim$500--2,000
tokens) and the clearest reward signal: did the chosen tactic lead to
a high-scoring speech? Group~B (skeleton + evidence) is grouped because
skeleton structure and evidence selection are tightly coupled---a skeleton
designed around unavailable evidence fails at generation time. Group~C
(speech generation) has the longest outputs and the most subjective
evaluation, requiring more epochs to converge. Group~D (CX dialogue)
is isolated because dialogue has fundamentally different dynamics than
monologue (see Appendix~\ref{app:failure_modes} for the retry
degradation phenomenon specific to CX).

\subsection{SFT-GRPO Interleaving}

Each group follows a two-stage training protocol within each iteration:

\begin{enumerate}
\item \textbf{SFT on Opus golden samples.} Claude Opus generates
    high-quality reference outputs for each call type. The model is
    fine-tuned (SFT) on these samples to establish a strong behavioral
    prior, producing the SFT+ checkpoint.

\item \textbf{GRPO on preference pairs.} The SFT+ model generates
    $k{=}4$ trial outputs per call (with logprob capture). These trials
    are scored by the agentic judge (Appendix~\ref{app:agentic_judge}),
    producing preference pairs. GRPO then optimizes the model on these
    pairs with the group-specific hyperparameters.

\item \textbf{Merge to canonical.} The GRPO LoRA is merged into the
    canonical base model, which becomes the starting point for the next
    group.
\end{enumerate}

The sequential ordering (A $\rightarrow$ B $\rightarrow$ C $\rightarrow$ D)
is deliberate: tactic improvements (Group~A) propagate to skeleton
construction (Group~B), which propagates to speech generation (Group~C),
and finally dialogue (Group~D). Each group trains on top of the
improvements from all preceding groups within the same iteration.

\paragraph{Why SFT before GRPO.}
GRPO alone can only redistribute probability mass among behaviors the
model already exhibits. SFT on Opus golden samples introduces
\emph{new} high-quality behaviors (e.g., more sophisticated argument
structures, better evidence integration) that GRPO can then reinforce
through preference optimization. Without the SFT step, GRPO is
limited to choosing among the model's existing capabilities.

\subsection{Importance Sampling}

GRPO requires the log-probability of each response under the current
policy. In the offline setting, responses were generated by a
\emph{previous} policy checkpoint. We handle this via importance sampling:

\begin{enumerate}
\item \textbf{Logprob capture at generation time.} When the SFT+ model
    generates trial outputs, per-token log-probabilities are recorded
    alongside the text.

\item \textbf{Importance weight computation.} The GRPO loss uses
    $\pi_\theta(y|x) / \pi_{\text{ref}}(y|x)$ ratios, where
    $\pi_{\text{ref}}$ is the generating policy. Precomputed logprobs
    from step~1 serve as $\log \pi_{\text{ref}}$.

\item \textbf{Recomputation after merge.} After each group merge, the
    canonical model's policy has shifted. For subsequent groups, logprobs
    from earlier generations become stale. We recompute reference logprobs
    on a random 20\% subsample to estimate the KL divergence; if
    $D_\text{KL} > 0.5$, all logprobs for the affected group are
    recomputed before training.
\end{enumerate}

Stale logprobs were identified as a root cause of the Group~C regression
described in Appendix~\ref{app:experiment_details}: the importance
weights became extreme, producing unstable gradient updates.

\subsection{Advantage Normalization}

Standard GRPO normalizes advantages globally across the training batch.
This is problematic for debate because the difficulty of generating a
good response depends heavily on the strategic context: a tactic that
works well against one opponent strategy may fail against another.

We normalize advantages within \emph{context triplets}:
$(\text{our\_tactic}, \text{opponent\_tactic}, \text{opponent\_intent})$.
For each triplet, the advantage is:
\begin{equation}
A_i = r_i - \bar{r}_{\text{triplet}}
\end{equation}

\noindent where $\bar{r}_{\text{triplet}}$ is the mean reward across all
samples sharing the same context triplet. This ensures that a mediocre
response against a strong opponent tactic is not penalized as harshly as
the same response against a weak opponent tactic.

In practice, triplet normalization improved validation accuracy by 3--5
percentage points across all groups compared to global normalization,
with the largest improvement in Group~A (tactic selection), where
strategic context has the most direct influence on reward.

\subsection{Attention-Only LoRA}

During the Group~C regression investigation
(Appendix~\ref{app:experiment_details}), we discovered that
targeting all seven linear projection modules in the Qwen3-30B-A3B
architecture (q\_proj, k\_proj, v\_proj, o\_proj, gate\_proj, up\_proj,
down\_proj) produced 3.37 billion trainable parameters. This large
parameter count, combined with stale offline logprobs, caused
catastrophic divergence.

Restricting LoRA to attention projections only (q\_proj, k\_proj,
v\_proj, o\_proj) reduced trainable parameters to \textbf{53 million}---a
$\mathbf{64\times}$ reduction. Critically, this produced \textbf{no measurable
quality loss}: validation accuracy on Group~C was 72.3\% with
attention-only LoRA versus 72.1\% with the full 7-module LoRA (before
regression), while completely eliminating the instability.

\begin{table}[h]
\centering
\small
\begin{tabular}{lccc}
\toprule
\textbf{LoRA Config} & \textbf{Params} & \textbf{Val Acc} & \textbf{Stable?} \\
\midrule
Full (7 modules) & 3.37B & 72.1\% & No (regression at epoch 3) \\
Attention-only (4 modules) & 53M & 72.3\% & Yes (all 8 epochs) \\
\bottomrule
\end{tabular}
\caption{Comparison of LoRA configurations for Group~C GRPO training.
Attention-only LoRA achieves equivalent accuracy with $64\times$ fewer
parameters and complete training stability.}
\label{tab:lora_comparison}
\end{table}

We hypothesize that for Mixture-of-Experts architectures with offline
GRPO, the MLP expert layers (gate\_proj, up\_proj, down\_proj) are
particularly sensitive to importance weight noise because each expert
is activated sparsely---small perturbations to expert gating can
cascade into large output changes. Attention-only LoRA avoids this
failure mode entirely.

\subsection{Per-Group Training Dynamics}

\paragraph{Group A (Tactic Selection).}
Converges fastest (3 epochs to plateau). The primary learning signal is
tactic-outcome correlation: the model learns which macro-tactics succeed
against which opponent strategies. Validation accuracy reaches 81.0\%,
the highest among speech groups, reflecting the relatively constrained
output space.

\paragraph{Group B (Skeleton + Evidence).}
Requires the longest max token length (16,000) due to evidence card
formatting. The evidence verification penalty (Appendix~\ref{app:method_details})
provides a strong shaping signal: fabrication and hallucination receive
$2\times$ GRPO weight, driving the model toward verified citations early
in training.

\paragraph{Group C (Speech Generation).}
The most challenging group, requiring 8 epochs to converge. Early epochs
show improvement in argument structure; later epochs improve rhetorical
quality. The epsilon clipping parameter ($\varepsilon = 0.15$) is critical:
values above 0.20 risk the regression described in
Appendix~\ref{app:experiment_details}.

\paragraph{Group D (CX Dialogue).}
Achieves the highest validation accuracy (87.5\%) but the least
improvement over the SFT baseline, consistent with the finding in
Appendix~\ref{app:experiment_details} that CX dialogue resists
preference optimization. We attribute this to the interactive nature
of dialogue: CX quality depends as much on the opponent's responses
as on the model's own outputs, making the reward signal noisier.

\subsection{DeepSpeed ZeRO-3 Checkpoint Fix}

Training uses DeepSpeed ZeRO-3 across 8 GPUs, which shards model
parameters, gradients, and optimizer states across ranks. A subtle
issue arises at checkpoint time: na\"ively calling
\texttt{model.save\_pretrained()} saves only the local shard, producing
a corrupted checkpoint. The fix is to explicitly gather the full state
dictionary before saving:

\begin{enumerate}
\item Call \texttt{deepspeed.zero.GatheredParameters()} to reconstruct
    the full parameter tensors on rank~0.
\item Save the gathered state dict from rank~0 only.
\item Resume sharded training after the save completes.
\end{enumerate}

This issue is not specific to our pipeline but affects any ZeRO-3
training with LoRA adapters. We document it here because it caused
silent checkpoint corruption during early Group~B training runs,
producing models that loaded without error but generated degenerate
outputs.
